\documentclass[11pt]{article}

% Packages pour les marges
\usepackage[
    top=2cm,
    bottom=2cm,
    left=2cm,
    right=2cm
]{geometry}

% Police sans serif
% \usepackage{helvet}
% \renewcommand{\familydefault}{\sfdefault}

% Packages existants
\usepackage[french]{babel}
\usepackage[utf8]{inputenc}
\usepackage[T1]{fontenc}
\usepackage{amsmath}
\usepackage{amsfonts}
\usepackage{amssymb}
\usepackage[version=4]{mhchem}
\usepackage{stmaryrd}
\usepackage{listings}

\usepackage{enumitem}
\usepackage{ifthen}
\usepackage{graphicx}
\usepackage{verbatim}
\usepackage{url}
\usepackage{mdframed}
\usepackage[sf]{titlesec}
\titleformat{\section}{\sffamily\large\bfseries}{\thesection}{1em}{}
% Définition de la variable pour afficher les corrections
\newboolean{showSolutions}
% Décommentez la ligne suivante pour afficher les solutions
\input \jobname.adr
% \input \jobname.adr


\title{TD1 - Probabilités discrètes}

\author{}
\date{}

\newenvironment{solution}
    {\par\vspace{0.5em}\begin{mdframed}[linewidth=0.5pt]\par}
    {\end{mdframed}\par\vspace{0.5em}}

\begin{document}
\sffamily
\maketitle

\section*{1 - Manipulation d'événements}
On lance une pièce une infinité de fois. Pour $n \geq 1$, on note $A_n$ l'événement "le $n$-ème lancer amène pile".
\begin{enumerate}
    \item Comment écrire à partir des $A_i$ les évenements : 
    \begin{itemize}
        \item Les 3 premiers lancers donnent pile
        \item Parmi les 2 premiers lancers, le premier est pile et le second est face. 
        \item Parmi les 2 premiers lancers, 1 est pile et 1 face. 
    \end{itemize}

    \ifthenelse{\boolean{showSolutions}}{
        \begin{solution}
    \begin{itemize}
        \item $A_1 \cap A_2 \cap A_3$
        \item $A_1 \cap \bar{A_2}$
        \item $(A_1\cap \bar{A_2}) \cup (\bar{A_1}\cap A_2)$
    \end{itemize}
\end{solution}
    }{}
    \item Décrire par une phrase les événements suivants:
    $$
B=\bigcap_{i=5}^{+\infty} A_i, C=\bigcup_{i=5}^{+\infty} A_i
$$
    \ifthenelse{\boolean{showSolutions}}{
        \begin{solution}
    $B$ signifie qu'à partir du cinquième lancer, on n'obtient que des piles. $C$ signifie qu'un lancer au moins à partir du cinquième lancer donne pile. On peut décrire $D$ par la formule
        \end{solution}
    }{}


\item  Écrire à l'aide de $A_i$ l'événement $D$ "On n'obtient plus que des piles à partir d'un certain lancer".

\ifthenelse{\boolean{showSolutions}}{
    \begin{solution}

$$
D=\bigcup_{n=1}^{+\infty} \bigcap_{i \geq n} A_i\,.
$$

    \end{solution}
}{}

\item La pièce est équilibrée, calculer la probabilité que parmi les deux premiers lancers, l'un des deux soit pile. 

\item Quelle est la probabilité que le deuxième lancer soit pile sachant que le premier lancer est pile ? 

\item Sachant qu'un des deux premiers lancers est pile, quelle est la probabilité que l'autre lancer soit pile ? 
\end{enumerate}
\vspace{2em}

\section*{2 - Mesure de Probabilité}
Soit $a \in] 0,1[$.

On appelle mesure de probabilité sur un ensemble dénombrable $X = \{x_1, x_2, \cdots\}$ une fonction $P: X \rightarrow [0,1]$ qui soit sous-additive et telle que la somme des poids est égale à $1$:
\[
\sum_{k\geq 1} P\big(\{x_k\}\big) = 1
\]
\begin{enumerate}

\item Déterminer $a$ tel que la fonction $P$, qui à tout $n \in \mathbb{N}$ associe $P(\{n\})=(1-a) a^n$ soit une mesure de probabilité sur $\mathbb{N}$. 
\ifthenelse{\boolean{showSolutions}}{
    \begin{solution}
 Il suffit de vérifier que, pour tout $\left.n \in \mathbb{N}, a^n(1-a) \in\right] 0,1[$ - ce qui est clair puisque $a \in] 0,1[$, donc $\left.\left.a^n \in\right] 0,1\right]$ ( $n$ peut être égal à 0 ) et $\left.1-a \in\right] 0,1[$ - et que

$$
\sum_{n=0}^{+\infty}(1-a) a^n=\frac{1-a}{1-a}=1
$$
    \end{solution}
}{}

\item On considère les deux événements $A=\{2 k: k \in \mathbb{N}\}$ et $B=\{2 k+1: k \in \mathbb{N}\}$. Calculer $P(A)$ et $P(B)$. Les événements $A$ et $B$ sont-ils incompatibles ? sont-ils indépendants ?
\end{enumerate} 

\ifthenelse{\boolean{showSolutions}}{
    \begin{solution}
On a

$$
P(A)=\sum_{k=0}^{+\infty}(1-a) a^{2 k}=\frac{1-a}{1-a^2}=\frac{1}{1+a}
$$


Comme $B=\bar{A}$, on a

$$
P(B)=1-P(A)=\frac{a}{1+a}
$$


Puisque $A \cap B=\varnothing$, les événements $A$ et $B$ sont incompatibles. Puisque $P(A \cap B)=P(\varnothing)=0$ et $P(A) P(B) \neq 0$, les événements $A$ et $B$ ne sont pas indépendants.
    \end{solution}
}{}
\vspace{2em}

\section*{3 - Indépendance d'événements}
\begin{enumerate}
\item Une urne contient $12$ boules numérotées de $1$ à $12$. On en tire une hasard, et on considère les événements

$$
\begin{aligned}
& A=\text { "tirage d'un nombre pair", } \\
& B=\text { "tirage d'un multiple de } 3 \text { ". }
\end{aligned}
$$


Les événements $A$ et $B$ sont-ils indépendants?

\ifthenelse{\boolean{showSolutions}}{
    \begin{solution}
        On a :

$$
\begin{gathered}
A=\{2,4,6,8,10,12\} \\
B=\{3,6,9,12\} \\
A \cap B=\{6,12\}
\end{gathered}
$$
On a donc $P(A)=1 / 2, P(B)=1 / 3$ et $P(A \cap B)=1 / 6=P(A) P(B)$. Les événements $A$ et $B$ sont indépendants.

\end{solution}
}{}

\item  Reprendre la question avec une urne contenant 13 boules.
\end{enumerate} 

\ifthenelse{\boolean{showSolutions}}{
    \begin{solution}
Les événements $A, B$ et $A \cap B$ s'écrivent encore exactement de la même façon. Mais cette fois, on a : $P(A)=6 / 13, P(B)=4 / 13$ et $P(A \cap B)=2 / 13 \neq 24 / 169$. Les événements $A$ et $B$ ne sont pas indépendants. C'est conforme à l'intuition. It n'y a plus la même répartition de boules paires et de boules impaires, et dans les multiples de 3 compris entre 1 et 13 , la répartition des nombres pairs et impairs est restée inchangée.
    \end{solution}
}{}
\vspace{2em}

\section*{4 - Formule des probabilités totales}
Dans une entreprise deux ateliers fabriquent les mêmes pièces. 
L'atelier 1 fabrique en une journée deux fois plus de pièces que l'atelier 2. 
Le pourcentage de pièces défectueuses est $3 \%$ pour l'atelier 1 et $4 \%$ pour l'atelier 2. 
On prélève une pièce au hasard dans l'ensemble de la production d'une journée. 

Déterminer (\textit{On pourra représenter un arbre de probabilités})
\begin{enumerate}
    \item la probabilité que cette pièce provienne de l'atelier 1 ;
    \ifthenelse{\boolean{showSolutions}}{
        \begin{solution}
            Notons $A$ l'événement " la pièce provient de l'atelier 1 ", $B$ l'événement " la pièce provient de l'atelier 2" et $D$ l'événement " la pièce est défectueuse".
    
            L'énoncé nous dit que les $2 / 3$ des pièces produites proviennent de l'atelier 1. Donc $P(A)=2 / 3$
        \end{solution}
    }{}

    \item  la probabilité que cette pièce provienne de l'atelier 1 et est défectueuse;
    \ifthenelse{\boolean{showSolutions}}{
        \begin{solution}
    On cherche $P(A \cap D)=P_A(D) P(A)=0,03 \times \frac{2}{3}=\frac{1}{50}$.
        \end{solution}
    }{}

    \item  la probabilité que cette pièce provienne de l'atelier 1 sachant qu'elle est défectueuse.
\end{enumerate}

\ifthenelse{\boolean{showSolutions}}{
    \begin{solution}
On démontre de même que $P(B \cap D)=\frac{1}{75}$ et donc que

$$
P(D)=P(A \cap D)+P(B \cap D)=\frac{1}{30}
$$


Ainsi, on a

$$
P_D(A)=\frac{P(A \cap D)}{P(D)}=\frac{3}{5}
$$

    \end{solution}
}{}
\vspace{2em}

\section*{5 - Loi de variables aléatoires}
Soit $X$ une variable aléatoire suivant une loi uniforme sur $\{1, \ldots, 20\}$. Déterminer la loi de $\lfloor\sqrt{X}\rfloor$.

Vous pourrez vérifier vos calculs en montrant que la loi de $\sqrt{X}$ est bien une mesure de probabilité, sur quel ensemble ?

Calculer l'espérance de $X$ et l'espérance de $\lfloor\sqrt{X}\rfloor$.

On donne la formule pour calculer l'espérance : si une variable aléatoire $X$ peut prendre les valeurs $x_1, x_2, \cdots$, alors
\[
E(X) =  \sum_{k\geq 1}x_k P(X=x_k)\,. 
\]
En français, l'espérance est la moyenne des valeurs que peut prendre $X$, pondérée par les probabilités que $X$ prenne ces valeurs. 

\ifthenelse{\boolean{showSolutions}}{
    \begin{solution}
        Pour simplifier les notations, posons $Y=\lfloor\sqrt{X}\rfloor$. Alors, $Y$ peut prendre les valeurs $1,2,3,4$. Plus précisément,

    \begin{itemize}

\item  si $X=1, X=2$ ou $X=3$, alors $Y=1$. En particulier,

$$
P(Y=1)=P(X=1)+P(X=2)+P(X=3)=\frac{3}{20}
$$

\item  si $X=4, X=5, X=6, X=7$ ou $X=8$, alors $Y=2$. En particulier,

$$
P(Y=2)=\frac{5}{20}
$$

\item si $X=9, X=10, X=11, X=12, X=13, X=14$ ou $X=15$, alors $Y=3$. En particulier,

$$
P(Y=3)=\frac{7}{20}
$$

\item si $X=16,17,18,19$ ou 20 , alors $Y=4$. En particulier,

$$
P(Y=4)=\frac{5}{20}
$$
    \end{itemize}

Pour calculer l'espérance, on utilise la formule : $X$ peut prendre toutes les valeurs entre $1$ et $20$, chacune avec probabilité $1/20$, l'espérance de $X$ est donc : 
\begin{align*}
E(X) &=  1 \times  \frac{1}{20} +  2 \times  \frac{1}{20}  +  3 \times  \frac{1}{20}  + \cdots +  20 \times  \frac{1}{20} \\
    & = \frac{1}{20} \sum_{k=1}^{20} k \\
    & = \frac{1}{20} \frac{20\times 21}{2} \\ 
    & = \frac{21}{2}\,.
\end{align*}

Calculons maintenant l'espérance de $Y$ : on multiplie les valeurs que peut prendre $Y$ par les probabilités qu'on a calculé : 
\[
E(Y) = 1 \times \frac{3}{20} + 2 \times \frac{5}{20} + 3\times \frac{7}{20} + 4 \times \frac{5}{20}
\]


    \end{solution}
}{}

\vspace{2em}

\section*{6 - Un avant-goût de statistiques}
Un étang contient des brochets et des truites. On note $p$ la proportion de truites dans l'étang. On souhaite évaluer $p$, on prélève 20 poissons au hasard. On suppose que le nombre de poissons est suffisamment grand pour que ce prélèvement s'apparente à 20 tirages indépendants avec remise. 
On note $X$ le nombre de truites obtenues.
\begin{enumerate}
\item Quelle est la loi de $X$ ?
\item Le prélèvement a donné 8 truites. Pour quelle valeur de $p$ la quantité $P(X=8)$ est-elle maximale?
\item On réalise un deuxième prélèvement de $20$ poissons qui a donné $9$ truites, comment modifier votre estimation pour $p$ ?
\item On réalise $n$ prélèvements de $20$ poissons, comment obtenir une bonne estimation de $p$ ?
\end{enumerate}
\ifthenelse{\boolean{showSolutions}}{
    \begin{solution}
        1. On reconnait le principe d'un schéma de Bernoulli, avec 20 épreuves indépendantes et une probabilité $p$ de succès. $X$ suit donc une loi binomiale $\mathcal{B}(20, p)$.

        2. On a $P(X=8)=\binom{20}{8} p^8(1-p)^{12}$. On va étudier cette fonction de $p$ de sorte de trouver son maximum. Il suffit en réalité d'étudier, sur $[0,1]$, la fonction $g(p)=p^8(1-p)^{12} . g$ est dérivable et sa dérivée est

$$
\begin{aligned}
g^{\prime}(p) & =8 p^7(1-p)^{12}-12 p^8(1-p)^{11} \\
& =(8(1-p)-12 p) p^7(1-p)^{11} \\
& =(8-20 p) p^7(1-p)^{11}
\end{aligned}
$$


On en déduit que $g^{\prime}(p) \geq 0$ sur $[0,2 / 5]$ et $g^{\prime}(p) \leq 0$ sur $[2 / 5,1]$. La probabilité $P(X=8)$ est maximale pour $p=2 / 5$.
    \end{solution}
}{}
\vspace{2em}

\section*{7 - Loi de Poisson}
On rappelle que la loi de Poisson est une mesure de probabilité sur $\mathbb{N}$ telle que 
\[
P(n) = e^{-\lambda} \, \frac{\lambda^k}{k!}\,.
\]
\begin{enumerate}
\item[0.] Montrer que cette fonction est bien une mesure de probabilité. 

\vspace{1em}

Un insecte pond des oeufs. Le nombre d'oeufs pondus est une variable aléatoire $X$ suivant une loi de Poisson $\mathcal{P}(\lambda)$. Chaque oeuf a une probabilité $p$ d'éclore, indépendante des autres oeufs. Soit $Z$ le nombre d'oeufs qui ont éclos.
    \item Pour $(k, n) \in \mathbb{N}^2$, calculer $P(Z=k \mid X=n)$.
\item En déduire la loi de $Z$ ?
\item Quelle est l'espérance de $Z$ ?
\end{enumerate}

\ifthenelse{\boolean{showSolutions}}{
    \begin{solution}
        1. On doit déterminer le nombre d'oeufs éclos sachant que $n$ oeufs ont été pondus. On a alors une succession de $n$ épreuves de Bernoulli indépendantes (l'oeuf $i$ éclot) avec probabilité $p$ de réussite à chaque fois. On reconnait donc une loi binomiale et on a

$$
P(Z=k \mid X=n)=\binom{n}{k} p^k(1-p)^{n-k}
$$

si $0 \leq k \leq n, P(Z=k \mid X=n)=0$ sinon.
2. On va appliquer la formule des probabilités totales : pour $k \in \mathbb{N}$,

$$
\begin{aligned}
P(Z=k) & =\sum_{n \geq k} P(Z=k \mid X=n) P(X=n) \\
& =\sum_{n \geq k}\binom{n}{k} p^k(1-p)^{n-k} \frac{e^{-\lambda} \lambda^n}{n!} \\
& =\frac{e^{-\lambda} p^k \lambda^k}{k!} \sum_{n \geq k} \frac{(1-p)^{n-k} \lambda^{n-k}}{(n-k)!}
\end{aligned}
$$


On réindexe la somme en posant $u=n-k$. Il vient

$$
\begin{aligned}
P(Z=k) & =\frac{e^{-\lambda} p^k \lambda^k}{k!} \sum_{u \geq 0} \frac{(\lambda(1-p))^u}{u!} \\
& =\frac{e^{-\lambda} p^k \lambda^k}{k!} e^{\lambda-\lambda p} \\
& =\frac{e^{-\lambda p}(\lambda p)^k}{k!}
\end{aligned}
$$


Ainsi, $Z$ suit une loi de Poisson de paramètre $\lambda p$.
3. L'espérance de $Z$ est donc $\lambda p$.
    \end{solution}
}{}



\section*{8 - Densité continue}
On appelle densité une fonction $f \,: \,\mathbb{R} \rightarrow \mathbb{R}$ telle que 
\begin{itemize}
    \item[$\bullet$] $f$ est positive
    \item[$\bullet$] $f$ est continue (sauf en un nombre fini de points)
    \item[$\bullet$] $f$ est d'intégrale $1$. 
\end{itemize}

Les fonctions suivantes sont-elles des densités ? 
\begin{enumerate}
        \item $f_1(x) = \lambda \, e^{-\lambda x}$ définie sur $[0, \infty]$ avec $\lambda$ un paramètre positif. 
        \item $f_2(x) = \frac{1}{5}\, 1_{[-5, 5]}$ définie sur $\mathbb{R}$
        \item $f_3(x) = \begin{cases}\cos x & \text { si } x \in[0, \pi / 2] \\ 0 & \text { sinon }\end{cases}$
        \item $f_4(x) =\frac{1}{1+x^2}, x \in \mathbb{R}$
\end{enumerate}

\ifthenelse{\boolean{showSolutions}}
{
    \begin{solution}
        On vérifie d'abord que les fonctions données sont continues sauf en un nombre fini de points et positives sur $\mathbb{R}$. On doit donc vérifier que $\int_{-\infty}^{+\infty} f(t) d t$ converge et est égal à 1 pour déterminer s'il s'agit d'une densité.
        
        3. On a bien \[
        \int_0^{\pi / 2} \cos (x)=\sin (\pi / 2)-\sin (0)=1 
        \]
        $f_3$ définit bien une densité de probabilité. 
        % Notons $X_1$ une variable aléatoire de densité $f_3$. Sa fonction de répartition est alors donnée par $F_{X_1}(x)=0$ si $x<0, F_{X_1}(x)=\sin (x)$ si $x \in[0, \pi / 2]$ et $F_{X_1}(x)=1$ si $x \geq \pi / 2$. De plus, puisque $f_1$ est non-nulle seulement sur l'intervalle $[0, \pi / 2], X_1$ admet une espérance donnée par
        % $$
        % E\left(X_1\right)=\int_0^{\pi / 2} x \cos (x) d x=\frac{1}{2}(\pi-2)
        % $$
        
        % (intégrale qu'on peut calculer à l'aide d'une intégration par parties).
        
        4. On a
        
        $$
        \int_{-\infty}^{+\infty} \frac{d x}{1+x^2}=\lim _{x \rightarrow+\infty} \arctan (x)-\lim _{x \rightarrow-\infty} \arctan (x)=\pi \neq 1
        $$
        
        
        Ainsi, $f_4$ n'est pas une densité.
    \end{solution}
}{}

\section*{9 - Lois discrètes classiques}
On considère les lois discrètes suivantes :

\begin{enumerate}
    \item La loi binomiale $B(n, p)$ avec $n$ un entier positif et $p$ la probabilité de succès.
    \item La loi de Poisson $P(\lambda)$ avec $\lambda$ un paramètre positif.
    \item La loi géométrique $G(p)$ avec $p$ la probabilité de succès.
    \item La loi uniforme discrète $U(a, b)$ avec $a$ et $b$ des entiers tels que $a \leq b$.
\end{enumerate}

Pour chacune des lois précédentes, calculer l'espérance et la variance.

La variance mathématique d'une variable aléatoire discrète $X$ est la moyenne des écarts à la moyenne : 
\[
\mathop{Var}(X) = E\Big((X-E(X))^2\Big)
\]

En développant, on peut l'écrire aussi comme : 
\[
\mathop{Var}(X) = E(X^2) - E(X)^2
\]


Si $X$ prend des valeurs $x_i$, on a aussi :
\[
\mathop{Var}(X) = \sum_{i} \big(x_i - E(X)\big)^2 \cdot P(X=x_i)
\]


\ifthenelse{\boolean{showSolutions}}
{
    \begin{solution}
        On donne seulement les résultats, envoyez-moi vos calculs si vous ne tombez pas sur les mêmes.
        \begin{enumerate}
            \item Pour la loi binomiale $B(n, p)$, l'espérance est $E(X) = np$ et la variance est $Var(X) = np(1-p)$.
            \item Pour la loi de Poisson $P(\lambda)$, l'espérance est $E(X) = \lambda$ et la variance est $Var(X) = \lambda$.
            \item Pour la loi géométrique $G(p)$, l'espérance est $E(X) = \frac{1}{p}$ et la variance est $Var(X) = \frac{1-p}{p^2}$.
            \item Pour la loi uniforme discrète $U(a, b)$, l'espérance est $E(X) = \frac{a+b}{2}$ et la variance est $Var(X) = \frac{(b-a+1)^2-1}{12}$.
        \end{enumerate}
    \end{solution}
}{}

\end{document}